\chapter{Profetas do algoritmo}
\noindent\textit{Vinheta.} Você abre o celular ``só por cinco minutos''. Uma sequência perfeita de vídeos curtos prende seu olhar. Quando percebe, perdeu quarenta. Não foi acaso. Foi \textbf{curadoria algorítmica}: um profeta moderno que anuncia, dia e noite, o que você ``precisa'' ver.

\section{Mapa bíblico}
\begin{itemize}[leftmargin=*,noitemsep]
  \item \textbf{Falsos profetas e sinais} (Mt~24:11,24): falas sedutoras que desviam o coração.
  \item \textbf{Coceira nos ouvidos} (2~Tm~4:3--4): preferir narrativas que confortam desejos.
  \item \textbf{Cativo a Cristo, não à tela} (2~Co~10:5): levar todo pensamento cativo à obediência de Cristo.
  \item \textbf{Atenção como mordomia} (Ef~5:15--16): remir o tempo, porque os dias são maus.
\end{itemize}

\section{Como funciona a recomendação (sem jargão)}
Algoritmos aprendem com seus \textit{sinais de engajamento}: tempo de tela, curtidas, comentários, pausas, replays. Eles combinam dados seus com dados de milhões para prever o que mais prolonga sua permanência. \textbf{Objetivo}: maximizar atenção. \textbf{Efeito colateral}: reforço de vieses e \textbf{bolhas} que parecem profecias personalizadas.

\section{Quatro dinâmicas do ``profeta'' algorítmico}
\begin{itemize}[leftmargin=*,noitemsep]
  \item \textbf{Atração}: entrega conteúdos emocionalmente carregados (\emph{medo}, \emph{ira}, \emph{vaidade}).
  \item \textbf{Aderência}: recompensa você por voltar; cria \emph{rituais} de checagem.
  \item \textbf{Aderência tribal}: conecta você a grupos com linguagem interna (``nós vs. eles'').
  \item \textbf{Autoridade emergente}: o \emph{feed} vira referência: ``se está em alta, deve ser verdade''.
\end{itemize}

\begin{nicbox}
\textbf{NIC aplicado}\\
\textbf{Narrativa}: a plataforma narra o mundo a partir do seu engajamento, não da realidade.\\
\textbf{Identidade}: você é reduzido a um \emph{perfil preditivo}.\\
\textbf{Controle}: seu comportamento é \emph{nudgeado} (empurrado) em microdecisões diárias.
\end{nicbox}

\section{Mini-exegese: o poder da \textit{coceira} (2~Tm~4:3)}
Paulo descreve uma \textbf{demanda por narrativas} que se ajustem a desejos. Os ``profetas do algoritmo'' entregam precisamente isso: uma liturgia de confirmação. O antídoto continua sendo o mesmo: \textbf{Palavra saudável}, \textbf{sã doutrina}, \textbf{vida comunitária} e \textbf{sobriedade}.

\section{Ferramenta --- Higiene informacional em 7 passos}
\begin{checklist}
\begin{itemize}[leftmargin=*,noitemsep]
  \item \textbf{Limpeza de fonte}: siga menos, siga melhor. Desative notificações push não essenciais.
  \item \textbf{Janelas de consumo}: horários fixos (p.\,ex., 12h30 e 19h) --- fora disso, o app fica \emph{off}.
  \item \textbf{Tempo de tela}: limite por app (p.\,ex., 20--30 min/dia); senha administrada por alguém de confiança.
  \item \textbf{Variedade intencional}: insira fontes que discordam de você (2:1 entre concordância e dissonância).
  \item \textbf{Ritual 4P} (cap.~3): Provedor, Prova, Palavra, Propósito antes de compartilhar.
  \item \textbf{Desintoxicação semanal}: 12 horas sem \emph{feed}; use tempo para conversas e leitura bíblica.
  \item \textbf{Comunhão} acima de consumo: substitua 1 hora de rolagem por 1 encontro real.
\end{itemize}
\end{checklist}

\section{Ferramenta --- Plano 3 camadas (pessoa, família, igreja)}
\begin{checklist}[title={Protocolos práticos},breakable]
\begin{itemize}[leftmargin=*,noitemsep]
  \item \textbf{Pessoal}: escolha 2 horários fixos de ver redes; \emph{mute} palavras-gatilho; desinstale um app por 30 dias.
  \item \textbf{Família}: aparelho carregando fora do quarto; 1 refeição/dia sem tela; noite de checagem 4P (quinzenal).
  \item \textbf{Igreja}: encoraje \emph{estudos de caso} mensais (uma fake viral dissecada); política de comunicação clara para crises.
\end{itemize}
\end{checklist}

\section{Evite dois extremos}
\begin{itemize}[leftmargin=*,noitemsep]
  \item \textbf{Tecnofobia}: demonizar toda plataforma --- perde-se o campo missionário.
  \item \textbf{Tecnoeuforia}: abraçar toda novidade --- vira-se refém do \emph{hype}.
\end{itemize}
Empurre a tecnologia para o seu lugar: \textbf{serva}. Coloque Cristo no Dele: \textbf{Senhor}.

\section{Perguntas para grupos pequenos}
\begin{itemize}[leftmargin=*,noitemsep]
  \item Quais conteúdos mais sequestram sua atenção (medo, escândalo, vaidade)? Por quê?
  \item O que você \textit{pode} remover esta semana (notificação, app, horário gatilho)?
  \item Que política de crise sua igreja precisa escrever antes da próxima tempestade de boatos?
\end{itemize}

\bigskip
\noindent\textit{Próximo capítulo: Identidade programável --- IDs digitais, biometria e reputação.}
